\documentclass[10pt,letterpaper]{article}
%DIF LATEXDIFF DIFFERENCE FILE
%DIF DEL original-submission.tex   Wed Nov 18 12:23:56 2026
%DIF ADD disaster-tips.tex         Wed Nov 18 13:42:27 2026
\usepackage[top=0.85in,left=2.75in,footskip=0.75in]{geometry}

% amsmath and amssymb packages, useful for mathematical formulas and symbols
\usepackage{amsmath,amssymb}

% Use adjustwidth environment to exceed column width (see example table in text)
\usepackage{changepage}

% textcomp package and marvosym package for additional characters
\usepackage{textcomp,marvosym}

% cite package, to clean up citations in the main text. Do not remove.
\usepackage{cite}

% Use nameref to cite supporting information files (see Supporting Information section for more info)
\usepackage{nameref,hyperref}

% line numbers
\usepackage[right]{lineno}

% ligatures disabled
\usepackage[nopatch=eqnum, expansion=false]{microtype}
\DisableLigatures[f]{encoding = *, family = * }

% color can be used to apply background shading to table cells only
\usepackage[table]{xcolor}

% array package and thick rules for tables
\usepackage{array}

% create "+" rule type for thick vertical lines
\newcolumntype{+}{!{\vrule width 2pt}}

% create \thickcline for thick horizontal lines of variable length
\newlength\savedwidth
\newcommand\thickcline[1]{%
  \noalign{\global\savedwidth\arrayrulewidth\global\arrayrulewidth 2pt}%
  \cline{#1}%
  \noalign{\vskip\arrayrulewidth}%
  \noalign{\global\arrayrulewidth\savedwidth}%
}

% \thickhline command for thick horizontal lines that span the table
\newcommand\thickhline{\noalign{\global\savedwidth\arrayrulewidth\global\arrayrulewidth 2pt}%
\hline
\noalign{\global\arrayrulewidth\savedwidth}}


% Remove comment for double spacing
%\usepackage{setspace} 
%\doublespacing

% Text layout
\raggedright
\setlength{\parindent}{0.5cm}
\textwidth 5.25in
\textheight 8.75in

% Bold the 'Figure #' in the caption and separate it from the title/caption with a period
% Captions will be left justified
\usepackage[aboveskip=1pt,labelfont=bf,labelsep=period,justification=raggedright,singlelinecheck=off]{caption}
\renewcommand{\figurename}{Fig}

% Please use the included `plos2025.bst` as your BibTeX style
\bibliographystyle{plos2025}

% Remove brackets from numbering in List of References
\makeatletter
\renewcommand{\@biblabel}[1]{\quad#1.}
\makeatother



% Header and Footer with logo
\usepackage{lastpage,fancyhdr,graphicx}
\usepackage{epstopdf}
%\pagestyle{myheadings}
\pagestyle{fancy}
\fancyhf{}
\setlength{\headheight}{27.023pt}
\lhead{\includegraphics[width=2.0in]{PLOS-submission.eps}}
\rfoot{\thepage/\pageref{LastPage}}
\renewcommand{\headrulewidth}{0pt}
\renewcommand{\footrule}{\hrule height 2pt \vspace{2mm}}
\fancyheadoffset[L]{2.25in}
\fancyfootoffset[L]{2.25in}
\lfoot{\today}

%% Include all user-defined macros below

\usepackage{callouts-box}
\newcommand{\withurl}[2]{{#1}\footnote{{\texttt{#2}}}}

%DIF PREAMBLE EXTENSION ADDED BY LATEXDIFF
%DIF UNDERLINE PREAMBLE %DIF PREAMBLE
\RequirePackage[normalem]{ulem} %DIF PREAMBLE
\RequirePackage{color}\definecolor{RED}{rgb}{1,0,0}\definecolor{BLUE}{rgb}{0,0,1} %DIF PREAMBLE
\providecommand{\DIFadd}[1]{{\protect\color{blue}\uwave{#1}}} %DIF PREAMBLE
\providecommand{\DIFdel}[1]{{\protect\color{red}\sout{#1}}} %DIF PREAMBLE
%DIF SAFE PREAMBLE %DIF PREAMBLE
\providecommand{\DIFaddbegin}{} %DIF PREAMBLE
\providecommand{\DIFaddend}{} %DIF PREAMBLE
\providecommand{\DIFdelbegin}{} %DIF PREAMBLE
\providecommand{\DIFdelend}{} %DIF PREAMBLE
\providecommand{\DIFmodbegin}{} %DIF PREAMBLE
\providecommand{\DIFmodend}{} %DIF PREAMBLE
%DIF FLOATSAFE PREAMBLE %DIF PREAMBLE
\providecommand{\DIFaddFL}[1]{\DIFadd{#1}} %DIF PREAMBLE
\providecommand{\DIFdelFL}[1]{\DIFdel{#1}} %DIF PREAMBLE
\providecommand{\DIFaddbeginFL}{} %DIF PREAMBLE
\providecommand{\DIFaddendFL}{} %DIF PREAMBLE
\providecommand{\DIFdelbeginFL}{} %DIF PREAMBLE
\providecommand{\DIFdelendFL}{} %DIF PREAMBLE
%DIF COLORLISTINGS PREAMBLE %DIF PREAMBLE
\RequirePackage{listings} %DIF PREAMBLE
\RequirePackage{color} %DIF PREAMBLE
\lstdefinelanguage{DIFcode}{ %DIF PREAMBLE
%DIF DIFCODE_UNDERLINE %DIF PREAMBLE
  moredelim=[il][\color{red}\sout]{\%DIF\ <\ }, %DIF PREAMBLE
  moredelim=[il][\color{blue}\uwave]{\%DIF\ >\ } %DIF PREAMBLE
} %DIF PREAMBLE
\lstdefinestyle{DIFverbatimstyle}{ %DIF PREAMBLE
	language=DIFcode, %DIF PREAMBLE
	basicstyle=\ttfamily, %DIF PREAMBLE
	columns=fullflexible, %DIF PREAMBLE
	keepspaces=true %DIF PREAMBLE
} %DIF PREAMBLE
\lstnewenvironment{DIFverbatim}[1][]{\lstset{style=DIFverbatimstyle}}{} %DIF PREAMBLE
\lstnewenvironment{DIFverbatim*}[1][]{\lstset{style=DIFverbatimstyle,showspaces=true}}{} %DIF PREAMBLE
\lstset{extendedchars=\true,inputencoding=utf8}

%DIF END PREAMBLE EXTENSION ADDED BY LATEXDIFF

\begin{document}
\vspace*{0.2in}

\begin{flushleft}
  {\Large
    \textbf\newline{Twelve Quick Tips for Managing IT Disasters in Small Research Software Teams}
  }
  \newline
  \\
  {Greg Wilson}\textsuperscript{1*}
  \\
  \bigskip
  \textbf{1} Third Bit, Toronto, Ontario, Canada\\
  * Corresponding author, gvwilson@third-bit.com.
\end{flushleft}

In 2025, the US government \DIFdelbegin \DIFdel{launched }\DIFdelend \DIFaddbegin \DIFadd{initiated }\DIFaddend an unprecedented series of \DIFdelbegin \DIFdel{attacks on }\DIFdelend \DIFaddbegin \DIFadd{cuts to }\DIFaddend its own
scientific research groups.  A year later \withurl{GitHub}{https://github.com/}
dropped below 90\% availability for the first time, while wildfires in Canada,
France, Spain, and elsewhere forced researchers from the homes and labs.  These
events and others have reminded us just how fragile research computing systems
can be, and that planning for disasters is one of the most effective ways to
\DIFdelbegin \DIFdel{prevent them}\DIFdelend \DIFaddbegin \DIFadd{mitigate their impacts}\DIFaddend .

This paper is a short guide to disaster planning and recovery for a small
research software team.  \DIFdelbegin \DIFdel{The tips assume you are doing everything yourself on
top of your regular job, and }\DIFdelend \DIFaddbegin \DIFadd{It assumes }\DIFaddend that you \emph{aren't} an experienced system
administrator\DIFdelbegin \DIFdel{. Some of the tips do require that kind of expertise, but }\DIFdelend \DIFaddbegin \DIFadd{, and anything you do to prepare for the worst is extra work on top
of ``normal'' work.  Note that }\DIFaddend most research institutions have research
computing groups, data librarians, and \DIFdelbegin \DIFdel{environmental health-and-safety offices whose entire job is to }\DIFdelend \DIFaddbegin \DIFadd{other specialists who can }\DIFaddend help with
exactly these problems. \DIFdelbegin \DIFdel{This paper tells you what ``done'' looks like; they can often
provide it}\DIFdelend \DIFaddbegin \DIFadd{(In some places this support comes from national
infrastructure instead.)  The ``Questions for Support'' boxes throughout tell
you what to ask them for}\DIFaddend .

\DIFaddbegin \begin{notebox}[Priorities]
  \DIFadd{If you are only able adopt a handful of these tips, please make a risk list
  and a written plan (Tips 1-2), follow the 3-2-1 backup rule (Tip 3), choose
  one out-of-band way to stay in touch (Tip 4), and turn on multi-factor
  authentication with a shared password manager (Tip 7).
}\end{notebox}

\DIFaddend \section*{Who You Are}

We assume that your team fits one or more of these profiles:

\begin{description}

\item[Providing a simple online service]
  Your research group runs \DIFdelbegin \DIFdel{a }%DIFDELCMD < \withurl{Shiny}{https://shiny.posit.co/} %%%
\DIFdel{app
  or a }%DIFDELCMD < \withurl{Streamlit}{https://streamlit.io/} %%%
\DIFdel{dashboard }\DIFdelend \DIFaddbegin \DIFadd{an interactive web site or web-based application
  }\DIFaddend so that collaborators can query data.
  Disasters scenarios include:
  \begin{itemize}
  \item Your cloud provider suspends your account.
  \item The graduate student who deployed the service leaves and no one else has credentials.
  \item Your service depends on a database that another group maintains,
    and that group shuts down on short notice.
  \end{itemize}

\item[Publishing a software package]
  Your group releases software to \DIFdelbegin \DIFdel{public registries }\DIFdelend \DIFaddbegin \DIFadd{a public registry }\DIFaddend like
  \withurl{PyPI}{https://pypi.org/} or \withurl{CRAN}{https://cran.r-project.org/}\DIFaddbegin \DIFadd{,
  to }\withurl{GitHub}{https://github.com/}\DIFadd{, or to an institutional download site}\DIFaddend .
  Scenarios include:
  \begin{itemize}
  \item \DIFaddbegin \DIFadd{A released package stops working after an operating system update, and
    the person who wrote it has left.
  }\item \DIFaddend Your PyPI token expires and the person who created it has left.
  \item GitHub suspends your account.
  \item Your institutional web hosting is turned off.
  \end{itemize}
  \DIFaddbegin \DIFadd{Note that once a package is published, the registry copy itself is usually not
  a single point of failure (CRAN, for example, has multiple mirrors).  The
  risks are the accounts and credentials around the release.
  }\DIFaddend 

\item[Managing data]
  Your group produces datasets in the lab or in the field\DIFdelbegin \DIFdel{that }\DIFdelend \DIFaddbegin \DIFadd{, whether or not they
  }\DIFaddend are shared with collaborators \DIFdelbegin \DIFdel{and }\DIFdelend \DIFaddbegin \DIFadd{or }\DIFaddend deposited in repositories.  Scenarios include:
  \begin{itemize}
  \item A lab server's hard drive fails.
  \item A field laptop is stolen and the data on it has not been backed up.
  \item Your data repository announces it is shutting down.
  \end{itemize}

\end{description}

These categories overlap: \DIFdelbegin \DIFdel{for example, }\DIFdelend an ecology lab may publish an R package, run a Shiny
dashboard for collaborators, and curate a ten-year field dataset.  \DIFaddbegin \DIFadd{At the other
extreme, a team of one or two people who only write analysis code and publish a
package may find all twelve tips disproportionate; such teams should focus on
the essential core named in the introduction and ignore, defer, or delegate the
rest.
}\DIFaddend 

\section*{Tip 1: Know your risks.}

Start by making a point-form list of the services, data sets, and physical
assets your team depends on.  For a \DIFdelbegin \DIFdel{small team}\DIFdelend \DIFaddbegin \DIFadd{simple service}\DIFaddend , this will take about an hour
the first time through, and 15-30 minutes per quarter for review.  \DIFdelbegin \DIFdel{For our ecology lab , this }\DIFdelend \DIFaddbegin \DIFadd{Using an
ecology lab as an example once again, this list }\DIFaddend includes where the R package is
published, where the Shiny dashboard is hosted, where the ten-year field dataset
lives, and the laptops the team works on. For each item, answer two questions:

\begin{enumerate}

\item
  How long can you \DIFdelbegin \DIFdel{be }\DIFdelend \DIFaddbegin \DIFadd{tolerate being }\DIFaddend without it before \DIFdelbegin \DIFdel{work stops dead}\DIFdelend \DIFaddbegin \DIFadd{your work is seriously
  disrupted}\DIFaddend ? This is your \emph{Recovery Time Objective} (RTO).

\item
  How much data can you afford to lose: the last hour's work, the last
  day's, a week's? This is your \emph{Recovery Point Objective} (RPO).

\end{enumerate}

A Google Doc or a text file in a shared drive or Git repository is the right
tool for creating this, but as discussed below, make sure you keep a copy
somewhere you can access when online services are down.  Once you have this
list, identify the single points of failure within your control, such as the
person who knows the deployment process or the credit card that pays for cloud
services. For each, ask what happens if it disappears. \DIFaddbegin \DIFadd{This is an ongoing
assessment, not a one-off: review the list quarterly and whenever your tools,
data, or people change.
}\DIFaddend 

\begin{notebox}[Questions for Support]
  If you have access to \DIFdelbegin \DIFdel{research }\DIFdelend \DIFaddbegin \DIFadd{IT }\DIFaddend support teams of the kind described in the introduction,
  \DIFdelbegin \DIFdel{you may want to ask themthese questions.
  }\DIFdelend \DIFaddbegin \DIFadd{ask them:
  }\DIFaddend \begin{itemize}
  \item Can you help us estimate realistic recovery time and recovery point targets?
  \item Who else shares the systems we depend on, and how would their failure cascade to us?
  \end{itemize}
\end{notebox}

\section*{Tip 2: Make a plan.}

Once you understand your risks, write a plan in a single shared document that
everyone on the team can find \DIFdelbegin \DIFdel{in thirty seconds or less}\DIFdelend \DIFaddbegin \DIFadd{quickly}\DIFaddend . Store the plan in at least two places
that cannot fail simultaneously, such as the shared team drive \emph{plus} a
printed copy in a desk drawer \DIFaddbegin \DIFadd{at home }\DIFaddend or a PDF saved on every team member's
phone. Pin the link in your team chat, and if you have a physical office, tape a
printed copy to the fridge.

State the disaster declaration criteria in plain language, such as ``our
dashboard has been offline for more than a day'' or ``we have found a virus on a
server''.  Every team member is part of the recovery team, so make sure everyone
has read the plan, understands the criteria, and knows their first action if a
disaster is declared.

Your plan should include \DIFaddbegin \DIFadd{contact information for the people who manage key
accounts and infrastructure.  It should include }\DIFaddend a short checklist for each
scenario, written as a numbered series of specific actions. For example, for an
online service outage:

\begin{enumerate}
\item Send a message to \DIFdelbegin \DIFdel{the }\DIFdelend \DIFaddbegin \DIFadd{your out-of-band channel (for example, a
  }\DIFaddend \withurl{Signal}{https://signal.org/} group\DIFaddbegin \DIFadd{)}\DIFaddend .
\item Log into the database management console and click ``Restore from
  snapshot''.
\item Check that the dashboard loads and returns data correctly.
\end{enumerate}

\begin{notebox}
  If a step requires knowledge or access that only one person has, you have a
  \emph{lottery-factor} problem (also called a ``bus factor'' problem). Ask
  yourself, ``If Alice wins the lottery and moves to Bali tomorrow, can the rest
  of the team keep things running?''  This is the hardest part of planning: Alice
  probably doesn't recognize all the things only she knows how to do, and no one
  else does either because she has always taken care of it.
\end{notebox}

\begin{notebox}[Questions for Support]
  \begin{itemize}
  \item Will you review our plan and checklists and flag any steps that are
    wrong or missing?
  \item Where should we keep a copy so your staff can also find it during an
    incident?
  \item \DIFdelbegin \DIFdel{Who is our first point of contact on }\DIFdelend \DIFaddbegin \DIFadd{What is the way to reach }\DIFaddend your team when we declare a disaster\DIFaddbegin \DIFadd{, including
    when internal IT is down or we are locked out}\DIFaddend ?
  \end{itemize}
\end{notebox}

\section*{Tip 3: Back up everything.}

Follow the \emph{3-2-1 rule} \cite{krogh2009}: three copies of every critical
asset, on two different types of media, with one copy off-site. This applies to
data, software releases, electronic lab notebooks, and configuration
information.  For \DIFdelbegin \DIFdel{our }\DIFdelend \DIFaddbegin \DIFadd{an }\DIFaddend ecology lab, the three copies might be:

\begin{enumerate}
\item
  The working copy on a lab server or shared drive.
\item
  An automated backup. If your institution offers a research computing backup
  service, use it. If not, use a managed database (e.g., \withurl{Amazon
    RDS}{https://aws.amazon.com/rds/} or \withurl{Google Cloud
    SQL}{https://cloud.google.com/sql}) with automatic daily snapshots\DIFdelbegin \DIFdel{(no scripting required)}\DIFdelend \DIFaddbegin \DIFadd{, or take
  local snapshots on a database system or a NAS server whose filesystem supports
  them}\DIFaddend . For files, GUI tools like \withurl{Duplicati}{https://duplicati.com/} or
  your cloud provider's desktop sync app are sufficient.
\item
  An off-site \DIFdelbegin \DIFdel{deposit. For datasets, deposit snapshots }\DIFdelend \DIFaddbegin \DIFadd{copy. For raw data, keep a second backup in a different location
  or cloud account, separate from your working copy and your automated backup.
  Publishing is a related but distinct step: deposit the }\emph{\DIFadd{curated}} \DIFadd{version
  of a dataset }\DIFaddend in an institutional or domain repository
  (\withurl{Dataverse}{https://dataverse.org/},
  \withurl{Dryad}{https://datadryad.org/},
  \withurl{Zenodo}{https://zenodo.org/}, or \withurl{OSF}{https://osf.io/}) \DIFdelbegin \DIFdel{that
  provides a }\DIFdelend \DIFaddbegin \DIFadd{to
  get a citable }\DIFaddend DOI. For software, connect Zenodo or
  \withurl{Figshare}{https://figshare.com/} to your repository so every tagged
  release is automatically archived.
\DIFaddbegin \end{enumerate}

\DIFaddend For lab notebooks, scan or photograph key pages quarterly and save them with
your other backups.
\DIFdelbegin %DIFDELCMD < \end{enumerate}
%DIFDELCMD < %%%
\DIFdelend 

Back up your source code to at least two different software forges (e.g.,
\withurl{GitLab}{https://gitlab.com/} \DIFaddbegin \DIFadd{or
}\withurl{Codeberg}{https://codeberg.org/} \DIFaddend as well as GitHub) and turn on
automatic mirroring. A \texttt{git push} to \DIFdelbegin \DIFdel{a secondary remote }\DIFdelend \DIFaddbegin \DIFadd{the second forge }\DIFaddend is the most
cost-effective one-liner you will ever write, but automatic mirroring is even
better.

Back up your configuration as well, such as \DIFdelbegin \DIFdel{DNS records, }\DIFdelend environment variables, \DIFaddbegin \DIFadd{server
configuration, }\DIFaddend and pipeline definitions. These are small\DIFdelbegin \DIFdel{in size but catastrophic }\DIFdelend \DIFaddbegin \DIFadd{, but difficult }\DIFaddend to
reconstruct from memory. A single export once a week, stored alongside your
other backups, is probably sufficient. \DIFaddbegin \DIFadd{Store backups under a separate account
whose credentials are not kept on any production server, so a compromise of your
main account does not also reach your backups.
}\DIFaddend 

\DIFaddbegin \DIFadd{Describe your computational pipelines in a workflow description language (for
example, }\withurl{Nextflow}{https://www.nextflow.io/}\DIFadd{) and register them with
}\withurl{Dockstore}{https://dockstore.org/} \DIFadd{or a similar registry. A declarative
pipeline that is registered is one you can actually run again when you need it.
}

\DIFaddend Finally, \emph{test a full restore from backup} at least once a year: according
to one recent report, almost 20\% of backups fail to restore completely
\cite{backblaze2024}.  And remember that you may need to restore several times
in a real disaster, so check that your restore process is \emph{idempotent},
i.e., that you can run \DIFdelbegin \DIFdel{it and then run it again}\DIFdelend \DIFaddbegin \DIFadd{multiple times with the same result}\DIFaddend .  A recovery that
corrupts data the second time is as bad as no recovery at all.

\DIFdelbegin %DIFDELCMD < \begin{notebox}
%DIFDELCMD <   %%%
\DIFdel{If you can store backups on a separate cloud account with credentials that are
  not stored on any production server, do so. If an attacker compromises your main
  account, they should not get your backups as well.
}%DIFDELCMD < \end{notebox}
%DIFDELCMD < 

%DIFDELCMD < %%%
\DIFdelend \begin{notebox}[Questions for Support]
  \begin{itemize}
  \item What backup and off-site repository services do you provide?
  \item Can you test our restore process?
  \item Who manages our off-site deposits, and how do we reach them when we need
    a restore?
  \end{itemize}
\end{notebox}

\section*{Tip 4: Decide how to stay in touch.}

Choose a single communication channel that does not depend on your normal
infrastructure. For example, if you use \withurl{Slack}{https://slack.com/},
agree \emph{in advance} on \DIFaddbegin \DIFadd{an out-of-band channel such as }\DIFaddend a Signal group or a
phone tree for when Slack is unavailable.  Make sure the fallback is written in
your plan.

For lab and field staff, maintain a printed contact tree: \DIFaddbegin \DIFadd{who calls whom }\DIFaddend if the
server fails \DIFdelbegin \DIFdel{at
2:00 a.m., who calls whom}\DIFdelend \DIFaddbegin \DIFadd{outside working hours}\DIFaddend ? Post it on the lab fridge and save a photo
on every team member's phone.

Keep a printed (or phone-screenshot) contact list\DIFdelbegin \DIFdel{with every team member's
mobile number and a }\DIFdelend \DIFaddbegin \DIFadd{, but respect privacy: list
only the few key people who must be reachable, and include a phone-less fallback
such as a }\DIFaddend personal email address. Do not store \DIFdelbegin \DIFdel{this }\DIFdelend \DIFaddbegin \DIFadd{the list }\DIFaddend only in a shared drive
that may be affected by an outage.

Pre-write messages and store them with the plan and in your password manager's
shared vault:

\begin{enumerate}
\item
  For the team: ``This is a declared incident: meet in the Signal group.''
\item
  For users of your online service: ``We are investigating an outage and
  will update you within 60 minutes.''
\item
  For collaborators waiting on data: ``We have lost connectivity temporarily.
  Data collected so far is safe; expect a delay of [X] days.''
\item
  For afterward: ``Service is restored and the situation is under control.
  Here is what happened and what we are doing to prevent recurrence.''
\end{enumerate}

Designate one person as the communicator so that everyone else can focus on
fixing the problem. The communicator does not need to be technical: they need to
be calm and reliable. Silence erodes trust faster than bad news, so update
stakeholders at regular intervals (for example, every 60 minutes or at 10:00 am
and 3:00 pm) even if the update is ``still working on it''.

\begin{notebox}[Questions for Support]
  \begin{itemize}
  \item How can we reach your team during an outage?
  \item Is there an institutional status page or mailing list our updates should
    be mirrored to?
  \item Who on your team needs to receive our incident messages?
  \end{itemize}
\end{notebox}

\section*{Tip 5: Test the plan.}

The most revealing test for a small team is to hand the recovery checklist to
the newest team member, \DIFdelbegin \DIFdel{give them }\DIFdelend \DIFaddbegin \DIFadd{have someone else break }\DIFaddend a test copy \DIFdelbegin \DIFdel{that they can break without
consequences}\DIFdelend \DIFaddbegin \DIFadd{of the service}\DIFaddend , and
ask \DIFdelbegin \DIFdel{them }\DIFdelend \DIFaddbegin \DIFadd{the newest member }\DIFaddend to execute the plan without help. Every point where they
get stuck is a documentation or automation gap.

If no one has joined your team recently, gather everyone for an hour-long
walkthrough. Pick the thing that scares you most\DIFdelbegin \DIFdel{as }\DIFdelend \DIFaddbegin \DIFadd{, e.g., }\DIFaddend ``It's Tuesday morning
and all of Monday's data has disappeared from the database.'' Walk through the
checklist step by step and fix anything that is missing, wrong, or stale\DIFdelbegin \DIFdel{then
and there}\DIFdelend .

Keep track of how long these tests take and compare against your RTO. If
restoring the database takes four hours but your RTO is two hours, you must
either fix the former or adjust the latter.

\begin{notebox}[Questions for Support]
  \begin{itemize}
  \item Can your staff join our next drill so we can find gaps together?
  \item Can you provide a sandbox where we can practice for disasters without
    affecting production?
  \end{itemize}
\end{notebox}

\section*{Tip 6: Watch for trouble.}

Monitor the things that \DIFdelbegin \DIFdel{would get you out of bed at 3:00 a.m. }\DIFdelend \DIFaddbegin \DIFadd{matter most when they break.  (These are the items on
the risk list you made in Tip 1, especially the ones with the shortest RTO and
RPO.) }\DIFaddend For online services, \DIFdelbegin \DIFdel{free-tier }\DIFdelend \DIFaddbegin \DIFadd{monitor up status, load, and response time so that a
misbehaving bot or a denial-of-service attack shows up early. Free-tier
}\DIFaddend monitoring like \withurl{Uptime Robot}{https://uptimerobot.com/} and
\withurl{Healthchecks.io}{https://healthchecks.io/} require no code and are
infinitely better than nothing.

\DIFaddbegin \begin{notebox}
  \DIFadd{Note that uptime monitoring will not catch an active intrusion, i.e., a hacker
  breaking into your system. However, such a break-in is often accompanied by a
  flood of traffic to a previously-unknown external site.
}\end{notebox}

\DIFaddend Check that your dataset's DOI still resolves and that your package is still
installable from its registry. If your software has a ``downloads per month''
badge, a sudden drop to zero may be the first sign that your registry listing is
broken \DIFaddbegin \DIFadd{because a token expired, an account was suspended, or the package was
removed or deprecated}\DIFaddend .

Check your cloud bill at least once a month: a misconfigured resource can
generate a catastrophic bill even for a small team. Set budget alerts in your
cloud provider's console at 50\% and 90\% of your normal monthly spend, sent to
at least two people. Most cloud providers offer free-tier budget alerting
adequate for a small team.

\begin{notebox}[Questions for Support]
  \begin{itemize}
  \item What monitoring and alerts do you already run, and can we subscribe to them?
  \item Can you help us set up alerts and uptime checks for our accounts?
  \item Which early-warning metrics would tell us something is about to fail?
  \end{itemize}
\end{notebox}

\section*{Tip 7: Secure your accounts.}

Enabling multi-factor authentication (MFA) on every account is the
highest-impact security measure you can take \cite{smalls2021}. An authenticator
app or passkey is better than SMS-based MFA, though SMS is better than no MFA at
all.

Use a password manager like \withurl{1Password}{https://1password.com/} or
\withurl{Bitwarden}{https://bitwarden.com/} for the whole team\DIFdelbegin \DIFdel{so that you can
share credentials securely}\DIFdelend \DIFaddbegin \DIFadd{. A shared vault
synchronizes entries across everyone's devices, so a credential is available
wherever it is needed. Sharing does widen exposure, so share only what a role
requires and revoke access when someone leaves}\DIFaddend . Every shared credential goes in
the password manager: never send credentials over chat or email. Print the
password manager's recovery code and store several copies somewhere physically
safe: if you lose access to the password manager, you lose access to everything.

If you have a custom \DIFaddbegin \DIFadd{web }\DIFaddend domain, know who manages it: institutional IT, a
commercial registrar, or someone on the team. Make sure \DIFdelbegin \DIFdel{that person's contact information
is in the disaster plan, that }\DIFdelend auto-renewal is enabled
\DIFdelbegin \DIFdel{, }\DIFdelend and that at least two people can access the account. Set a quarterly calendar
reminder to verify that auto-renewal is still working.

\DIFaddbegin \DIFadd{MFA and password managers centralize your points of failure; used without backup
access they can lock you out of everything at once. Pair them with the recovery
paths above rather than treating them as a silver bullet.
}

\DIFaddend Make sure at least two people have critical administrative credentials, such as
the cloud root account, the domain registrar login, the payment method for
infrastructure, and the password-manager administrator account.  These
credentials should not be tied to a single phone number or email address.  For
physical access, maintain a list of who has keys to the lab, who knows the alarm
code, and who is authorized to move hardware.

Finally, maintain an \DIFdelbegin \DIFdel{offboarding checklist }\DIFdelend \DIFaddbegin \DIFadd{onboarding checklist alongside an offboarding one}\DIFaddend . When
someone \DIFaddbegin \DIFadd{joins, grant their access, add them to the password manager and the
contact tree, and have them read the plan and run the newest-member drill from
Tip 5. When someone }\DIFaddend leaves, revoke their digital access and collect their
physical keys as soon as you can. A disgruntled former team member with
lingering access is trouble you don't need.

\begin{notebox}[Questions for Support]
  \begin{itemize}
  \item Can you help us set up multi-factor authentication?
  \item Which of your team should have access to administrative credentials?
  \item How quickly does your offboarding process revoke access when someone
    leaves?
  \end{itemize}
\end{notebox}

\section*{Tip 8: Protect your assets.}

\DIFdelbegin \DIFdel{Make a simple, explicit policy about where things live, such as, ``patient data
is never copied to a personal laptop'' or ``database dumps are encrypted at
rest''.  If your policy will not fit comfortably on a single page in an 11-point
font, rewrite it until it does.
}%DIFDELCMD < 

%DIFDELCMD < %%%
\DIFdelend \begin{notebox}
  \DIFdelbegin \DIFdel{You }\DIFdelend \DIFaddbegin \DIFadd{If your institution's IT group owns your servers and laptops, several of the
  steps below become things to verify with or request from them rather than do
  yourself.  In particular, you }\DIFaddend may have obligations you are not aware of under
  HIPAA, GDPR, or your funder's data-management plan.  If you are handling
  sensitive data, \DIFdelbegin \DIFdel{you should
  therefore }\DIFdelend check whether your institution has an ethics board or
  data-protection office that offers free consultation.
\end{notebox}

\DIFaddbegin \DIFadd{Write a simple, explicit policy about where things live, such as, ``patient data
is never copied to a personal laptop'' or ``database dumps are encrypted at
rest''. Rules like the first cannot be left to good intentions alone: enforce
them technically, for example with a trusted research environment for protected
data.  Keep the policy short enough that people will actually read and remember
it.
}

\DIFaddend Keep laptops and servers patched: \DIFaddbegin \DIFadd{apply security updates promptly, and }\DIFaddend enable
automatic updates for operating systems, browsers, and \DIFdelbegin \DIFdel{any }\DIFdelend server-side software
\DIFaddbegin \DIFadd{where you can review and roll them back}\DIFaddend . The inconvenience of an unexpected
reboot is less than the inconvenience of ransomware.

For on-premises equipment like a lab server, consider buying a cheap
uninterruptible power supply (UPS) and configure it to trigger a clean shutdown
when the battery reaches 20\%.  Replace the UPS battery every three years: they
have a tendency to degrade silently.

\DIFaddbegin \DIFadd{Think about physical security too: if a critical server sits in a space many
people can enter, lock it away or restrict access to the room.
}

\DIFaddend Prepare a short, written response for ransomware and similar breaches. If
someone on the team is contacted by villains, their first action should be to
\DIFdelbegin \DIFdel{physically }\DIFdelend disconnect the affected machine from the \DIFdelbegin \DIFdel{network and }\DIFdelend \DIFaddbegin \DIFadd{network---physically, if it is in front
of you, or via its cloud console or security group, or by asking IT if it is
racked---and }\DIFaddend then call the designated incident lead. Do not try to negotiate or
pay without consulting legal counsel: most jurisdictions have regulations about
ransomware payments.

\DIFaddbegin \DIFadd{AI agents are the newest source of trouble. If your lab uses coding or analysis
agents that can read files, edit code, or run commands, treat them as you would
a new, overconfident team member: give them least-privilege access, log what
they do, and keep sensitive data and credentials out of the directories they can
reach. A poorly prompted agent can exfiltrate data, disclose confidential
information, or damage local files by accident.
}

\DIFaddend \begin{notebox}[Questions for Support]
  \begin{itemize}
  \item Who should we talk to about data protection?
  \item Can you confirm our laptops and servers have automatic patching enabled?
  \item If we face ransomware, who is our incident lead and which legal counsel
    do we call?
  \end{itemize}
\end{notebox}

\section*{Tip 9: Stabilize, then investigate.}

The most important rule of incident response is, ``Stabilize first, investigate
later''.  When something \emph{does} happen, assign clear roles for the duration
of the incident: one person leads the technical response, one person handles
communication using the pre-written messages from Tip 4, and everyone else does
what the leads tell them to.

Create an incident log \emph{as you work}: a shared Google Doc or a thread in
your out-of-band Signal group is good enough.  This log keeps the team in sync
and gives you a record for the post-incident review.  Timestamp every
significant action:

\begin{itemize}
\item 14:32: server outage alarm triggered
\item 14:35: verified that all database requests are returning 404
\item 14:40: began transition to backup server
\end{itemize}

Escalate early. If you are on a cloud provider's support plan, open a ticket the
moment you suspect the problem is on their side. You can cancel your request if
you fix the problem yourself, but you cannot get back the two hours you spent
wrestling with something outside your control.

Finally, conduct a short review within 48 hours of any incident. The goal is to
identify what in the system allowed the incident to happen, not who made a
mistake. Write down one concrete action item and assign it to one person with a
deadline. If you cannot find an outside moderator for the review, have the team
member \emph{least} involved in the incident lead it.

\begin{notebox}[Questions for Support]
  \begin{itemize}
  \item How do we escalate a ticket with our cloud provider(s)?
  \item Who on your team should join our incident response?
  \item Can you act as moderator for our post-incident review?
  \end{itemize}
\end{notebox}

\section*{Tip 10: Count the cost.}

You probably don't have a budget line item titled ``disasters'', but you still
have costs.  Digital costs include cloud backup storage, password-manager
subscription, domain and certificate renewals, managed-database fees, and
repository deposit charges.  More important, but harder to track, are the
personnel costs: time spent on quarterly plan reviews, annual backup-restore
tests and disaster drills, cross-training to reduce lottery factors, and the
hours lost to incident response\DIFdelbegin \DIFdel{instead of research}\DIFdelend .

Calculate what a day of downtime costs you, including things like lost
experiment time, missed paper deadlines, or collaborators who stop trusting
you. That number tells you whether it is worth spending \$50/month on a managed
database or \$200/year on a backup service.

\begin{notebox}[Questions for Support]
  \begin{itemize}
  \item Can you tell us what backup storage, a managed database, and monitoring
    ought to cost?
  \item Which of those can the institution cover, and which do we have to fund
    ourselves?
  \end{itemize}
\end{notebox}

\section*{Tip 11: Support your team during the crisis.}

The \DIFdelbegin \DIFdel{disasters discussed in the opening are not just logistics problems: each is
a shock to your team and to the peoplewho depend on them}\DIFdelend \DIFaddbegin \DIFadd{first ten tips are about systems and processes; these last two are about
people, because a disaster is as much a human event as a technical one}\DIFaddend .  Team
members may be scared, angry, or ashamed, particularly if they believe they
caused the problem.  Handling the human damage is therefore just as important as
handling the technical aspects.

During the acute phase of a disaster, apply the same five principles that the
World Health Organization recommends for crisis responders \cite{who2026}:

\begin{description}

\item[Promote safety.]
  Before asking for status updates, confirm that everyone is not in immediate
  personal danger.  Someone whose laptop was stolen may not be thinking clearly
  about data loss because they are still shaken from the theft itself.

\item[Calm the environment.]
  Do not add to the panic.  Use a steady voice in your out-of-band chat.  Remind
  the team that the plan exists and that they have rehearsed for this.  Panic is
  contagious, but so is calm.

\item[Build self-efficacy.]
  Remind people that they have handled difficult things before.  Give each person
  one concrete task from the checklist so that everyone feels useful rather than
  helpless.

\item[Foster connectedness.]
  Use your fallback communication channel to keep the team together.  Isolation
  amplifies distress, so do not let anyone absorb the news alone.

\item[Maintain hope.]
  In particular, do not speculate about worst cases before you have confirmed
  them.

\end{description}

Protect the person who may have caused the problem.  A team member who
accidentally dropped a table from a database or clicked a phishing link probably
feels worse than anyone else.  Assign them a recovery task, even a small one, so
they can contribute to the fix rather than sitting in the corner imagining their
career is over.  Blame can wait; recovery cannot.

Watch for exhaustion.  Disaster response is an adrenaline sprint followed by a
long slog.  Two people working four hours each will make fewer mistakes than one
person working eight hours straight, so rotate people out of the incident
response when the acute phase is over, even if they insist they are fine.  In
particular, the person who pulled the all-nighter to restore the database should
not also be the person writing the post-mortem the next morning.

\begin{notebox}[Questions for Support]
  \begin{itemize}
  \item What wellbeing or crisis support do you offer, and how do we access it?
  \item Are there any provisions for post-incident leave?
  \item How should crisis preparation and response be factored into annual
    performance reviews?
  \end{itemize}
\end{notebox}

\section*{Tip 12: Help your team recover after the disaster.}

\DIFdelbegin \DIFdel{In }\DIFdelend Tip 9 said to hold \DIFdelbegin \DIFdel{the }\DIFdelend \DIFaddbegin \DIFadd{a short }\DIFaddend review within 48 hours. That is for gathering facts
while memory is fresh; a second, more reflective conversation a week or two
later gives people time to process what happened emotionally.  The goal of this
second conversation is \DIFdelbegin \DIFdel{not to find root causes but to ask: }\DIFdelend \DIFaddbegin \DIFadd{to ask, }\DIFaddend ``How are we doing?''

If someone on the team caused the disaster, handle it carefully, and make sure
the rest of the team acts that way.  Ask them privately if they are OK.  If you
find yourself wanting to revisit the mistake in every team meeting, you are the
one who hasn't moved on.

Recognize that people recover at different speeds.  One team member may be
joking about it by Tuesday while another is still waking up anxious on Friday.
Both responses are normal.  Check in with each person individually a week after
the incident, and again a month later.  Again, the question is simply: ``How are
you doing with what happened?''

Watch for survivors' guilt.  For example, if your lab's data survived a hacking
intrusion but a neighbouring group lost everything, your team may not want to
celebrate their own recovery when others are still suffering.  Let them talk
about it. If possible, offer practical help to affected colleagues, such as
sharing backup space or simply bringing them coffee.

Finally, celebrate the recovery, not just the disaster. A short thank-you
message that says ``we handled this'' costs nothing and pays dividends the next
time something goes wrong.

\begin{notebox}[Questions for Support]
  \begin{itemize}
  \item What longer-term counselling or support can you offer staff in the weeks
    after the incident?
  \item Who else in the institution should we inform so they can support our
    recovery?
  \item Can you help us capture lessons learned so other teams benefit from what
    happened?
  \end{itemize}
\end{notebox}

\nocite{*}
\bibliography{disaster-tips}

\end{document}
